% part: intuitionistic-logic
% chapter: introduction
% section: bhk-interpretation

\documentclass[../../../include/open-logic-section]{subfiles}

\begin{document}

\olfileid{int}{int}{bhk}

\olsection{تفسیر براور--هایتینگ--کولموگروف}

\begin{editorial}
  اثبات اعتبار گزاره‌های شهودگرایانه با تفسیر BHK گیج‌کننده است؛ باید
  بهتر توضیح داده شود.
\end{editorial}

تفسیری سازنده و غیرصوری از رابط‌های شهودگرایانه وجود دارد که معمولاً
تفسیر براور--هایتینگ--کولموگروف نامیده می‌شود. این تفسیر از مفهوم
«ساخت» استفاده می‌کند که می‌توانید آن را اثباتی سازنده بدانید. (در
تفسیر BHK واژهٔ «اثبات» را به کار نمی‌بریم تا با مفهوم
!!a{derivation} در یک دستگاه صوریِ !!{derivation} خلط نشود.) تفسیر BHK
بر پایهٔ این مفهوم شهودی، معانی رابط‌های شهودگرایانه را توضیح می‌دهد.

\begin{enumerate}
\item فرض می‌کنیم می‌دانیم چه چیزی ساختِ یک گزارهٔ اتمی به شمار می‌آید.
\item ساختِ~$!A_1 \land !A_2$ جفت~$\tuple{M_1, M_2}$ است، که در آن
  $M_1$ ساختِ~$!A_1$ و~$M_2$ ساختِ~$!A_2$ است.
\item ساختِ~$!A_1 \lor !A_2$ جفت~$\tuple{s, M}$ است، که در آن یا
  $s$ برابر~$1$ و~$M$ ساختِ~$!A_1$ است، یا~$s$ برابر~$2$ و~$M$
  ساختِ~$!A_2$ است.
\item ساختِ~$!A \lif !B$ تابعی است که ساختی از~$!A$ را به ساختی
  از~$!B$ تبدیل می‌کند.
\item برای~$\lfalse$ (محال) هیچ ساختی وجود ندارد.
\item $\lnot !A$ هم‌معنای~$!A \lif \lfalse$ تعریف می‌شود. یعنی ساختِ
  $\lnot !A$ تابعی است که ساختی از~$!A$ را به ساختی از~$\lfalse$
  تبدیل می‌کند.
\end{enumerate}

\begin{ex}
برای نمونه~$\lnot \lfalse$ را در نظر بگیرید. ساختی برای آن تابعی است
که هر ساختِ~$\lfalse$ را به‌عنوان ورودی بگیرد و ساختی از~$\lfalse$ را
به‌عنوان خروجی بدهد. آشکارا تابع همانی~$\Id{}$ چنین ساختی است: با گرفتن
ساخت~$M$ از~$\lfalse$، رابطهٔ~$\Id{}(M) = M$ ساختی از~$\lfalse$
به دست می‌دهد.
\end{ex}

به‌طور کلی، $\lnot !A$ یعنی «ساختِ~$!A$ ناممکن است».

\begin{ex}
برای هر گزارهٔ~$!A$، $!A \lif \lnot \lnot !A$، یعنی
$!A \lif ((!A \lif \lfalse) \lif \lfalse)$، را ثابت کنیم. ساخت مورد
نظر باید تابع~$f$ای باشد که با گرفتن ساخت~$M$ از~$!A$، ساخت~$f(M)$ از
$(!A \lif \lfalse) \lif \lfalse$ را بازگرداند. $f$ ساختِ
$(!A \lif \lfalse) \lif \lfalse$ را چنین می‌سازد: باید تابع~$g$ای
تعریف کنیم که اگر ساخت~$h$ از~$!A \lif \lfalse$ را به‌عنوان ورودی
بگیرد، ساختی از~$\lfalse$ بیرون دهد. $g$ را چنین تعریف می‌کنیم: ورودی
$h$ را بر ساخت~$M$ از~$!A$ (که پیش‌تر دریافت کردیم) اعمال کنید. چون
خروجی~$h(M)$ از~$h$ ساختی از~$\lfalse$ است، اگر~$M$ ساختی از~$!A$
باشد، $f(M)(h) = h(M)$ ساختی از~$\lfalse$ است.
\end{ex}

\begin{ex}
ساختی برای~$\lnot(!A \land \lnot !A)$، یعنی
$(!A \land (!A \lif \lfalse)) \lif \lfalse$، بدهیم. این ساخت تابع
$f$ای است که با گرفتن ساخت~$M$ از~$!A \land (!A \lif \lfalse)$
به‌عنوان ورودی، ساختی از~$\lfalse$ به دست می‌دهد. ساختِ عطف
$!B_1 \land !B_2$ جفت~$\tuple{N_1, N_2}$ است که~$N_1$ ساختِ~$!B_1$
و~$N_2$ ساختِ~$!B_2$ است. می‌توانیم توابع~$p_1$ و~$p_2$ را تعریف
کنیم که از یک ساختِ~$!B_1 \land !B_2$ به‌ترتیب ساخت‌های~$!B_1$ و~$!B_2$
را بازیابی کنند:
\begin{align*}
  p_1(\tuple{N_1, N_2}) &= N_1\\
  p_2(\tuple{N_1, N_2}) &= N_2
\end{align*}
$f$ چنین عمل می‌کند: نخست~$p_1$ را بر ورودی~$M$ اعمال می‌کند. حاصل،
ساختی از~$!A$ است. سپس~$p_2$ را بر~$M$ اعمال می‌کند و ساختی از
$!A \lif \lfalse$ به دست می‌آورد. چنین ساختی خود تابع~$p_2(M)$ است
که اگر ساختی از~$!A$ را به‌عنوان ورودی بگیرد، ساختی از~$\lfalse$ به
دست می‌دهد. به بیان دیگر، اگر~$p_2(M)$ را بر~$p_1(M)$ اعمال کنیم،
ساختی از~$\lfalse$ می‌گیریم. پس می‌توانیم
$f(M) = p_2(M)(p_1(M))$ تعریف کنیم.
\end{ex}

\begin{ex}
ساختی برای~$((!A \land !B) \lif !C) \lif (!A \lif
(!B \lif !C))$ بدهیم؛ یعنی تابع~$f$ای که ساخت~$g$ از
$(!A \land !B) \lif !C$ را به ساختی از~$(!A \lif (!B \lif !C))$
تبدیل کند. ساخت~$g$ خود تابعی است (از ساخت‌های~$!A \land !B$ به
ساخت‌های~$!C$). خروجی~$f(g)$ نیز تابع~$h_g$ای است از ساخت‌های~$!A$
به توابعی از ساخت‌های~$!B$ به ساخت‌های~$!C$.

خب، این گیج‌کننده است. باید تابع معینی~$h_g$ بسازیم که خروجی~$f$ برای
ورودی~$g$ خواهد بود. ورودی~$h_g$ ساخت~$M$ از~$!A$ است. خروجی~$h_g(M)$
باید تابع~$k_M$ای از ساخت‌های~$N$ از~$!B$ به ساخت‌های~$!C$ باشد.
بگذارید~$k_{g,M}(N) = g(\tuple{M, N})$. به یاد آورید که
$\tuple{M, N}$ ساختی از~$!A \land !B$ است. پس~$k_{g,M}$ ساختی از
$!B \lif !C$ است: ساخت‌های~$N$ از~$!B$ را به ساخت‌هایی از~$!C$
می‌فرستد. اکنون بگذارید~$h_g(M) = k_{g,M}$. این تابع ساخت‌های~$M$ از
$!A$ را به ساخت‌های~$k_{g,M}$ از~$!B \lif !C$ می‌فرستد. حال بگذارید
$f(g) = h_g$. این تابع ساخت‌های~$g$ از~$(!A \land !B) \lif !C$ را به
ساخت‌هایی از~$!A \lif (!B \lif !C)$ می‌فرستد. سرانجام!{}
\end{ex}

گزارهٔ~$!A \lor \lnot !A$ قانون طرد شق ثالث نامیده می‌شود. می‌توانیم
آن را برای برخی~$!A$های معین (برای نمونه،
$\lfalse \lor \lnot \lfalse$) اثبات کنیم، اما نه به‌طور کلی. دلیل آن
این است که فصل شهودگرایانه ساختی از یکی از دو طرف فصل می‌طلبد، حال آنکه
گزاره‌هایی وجود دارند که در حال حاضر نه قابل اثبات‌اند و نه قابل رد
(مثلاً حدس گلدباخ). بااین‌حال، خود قانون طرد شق ثالث را نیز نمی‌توان رد
کرد؛ یعنی~$\lnot \lnot (!A \lor \lnot !A)$ برقرار است.

\begin{ex}
برای اثبات~$\lnot \lnot (!A \lor \lnot !A)$، به تابع~$f$ای نیاز داریم
که ساختی از~$\lnot(!A \lor \lnot !A)$، یعنی از
$(!A \lor (!A \lif \lfalse)) \lif \lfalse$، را به ساختی از~$\lfalse$
تبدیل کند. به بیان دیگر، تابع~$f$ای می‌خواهیم چنان‌که اگر~$g$ ساختی از
$\lnot(!A \lor \lnot !A)$ باشد، $f(g)$ ساختی از~$\lfalse$ باشد.

فرض کنید~$g$ ساختی از~$\lnot(!A \lor \lnot !A)$ باشد؛ یعنی تابعی که
ساختی از~$!A \lor \lnot !A$ را به ساختی از~$\lfalse$ تبدیل کند. ساختِ
$!A \lor \lnot !A$ جفت~$\tuple{s, M}$ است که در آن یا~$s = 1$ و~$M$
ساختی از~$!A$ است، یا~$s = 2$ و~$M$ ساختی از~$\lnot !A$. بگذارید~$h_1$
تابعی باشد که ساخت~$M_1$ از~$!A$ را به ساختی از~$!A \lor \lnot !A$
می‌فرستد: $M_1$ را به~$\tuple{1, M_1}$ می‌فرستد. و بگذارید~$h_2$ تابعی
باشد که ساخت~$M_2$ از~$\lnot !A$ را به ساختی از
$!A \lor \lnot !A$ می‌فرستد: $M_2$ را به~$\tuple{2, M_2}$ می‌فرستد.

بگذارید~$k$ برابر~$\comp{h_1}{g}$ باشد: این تابع اگر ساختی از~$!A$
بگیرد، ساختی از~$\lfalse$ بازمی‌گرداند؛ یعنی ساختی از
$!A \lif \lfalse$ یا~$\lnot !A$ است. اکنون بگذارید~$l$ برابر
$\comp{h_2}{g}$ باشد. این تابع با گرفتن ساختی از~$\lnot !A$، ساختی از
$\lfalse$ به دست می‌دهد. چون~$k$ ساختی از~$\lnot !A$ است، $l(k)$
ساختی از~$\lfalse$ است.

در مجموع، توضیح دادیم چگونه می‌توان ساخت~$g$ از
$\lnot(!A \lor \lnot !A)$ را به ساختی از~$\lfalse$ تبدیل کرد؛ یعنی
تابع~$f$ که ساخت~$g$ از~$\lnot(!A \lor \lnot !A)$ را به ساخت~$l(k)$
از~$\lfalse$ می‌فرستد، ساختی از~$\lnot\lnot(!A \lor \lnot !A)$ است.
\end{ex}

چنان‌که می‌بینید، استفاده از تفسیر BHK برای نشان‌دادن اعتبار
شهودگرایانهٔ !!{formula}ها به‌سرعت دشوار و گیج‌کننده می‌شود. خوشبختانه
دستگاه‌های !!{derivation} بهتر و تفسیرهای معناشناختی دقیق‌تری برای منطق
شهودگرایانه وجود دارند.
\end{document}
